\documentclass[12pt,a4paper]{article}

% ── Packages ────────────────────────────────────────────────────────────────
\usepackage[a4paper, margin=2.5cm]{geometry}
\usepackage{amsmath, amssymb, amsfonts}
\usepackage{graphicx}
\usepackage{booktabs}
\usepackage{array}
\usepackage{longtable}
\usepackage{multirow}
\usepackage{tabularx}
\usepackage{float}
\usepackage{hyperref}
\usepackage{setspace}
\usepackage{titlesec}
\usepackage{fancyhdr}
\usepackage{caption}
\usepackage{subcaption}
\usepackage{xcolor}
\usepackage{parskip}
\usepackage{enumitem}
\usepackage[T1]{fontenc}
\usepackage[utf8]{inputenc}
\usepackage{bm}
\usepackage{mathtools}
\usepackage{pgfplots}
\pgfplotsset{compat=1.18}
\usepackage{tikz}
\usetikzlibrary{arrows.meta, positioning}

% ── Hyperref setup ───────────────────────────────────────────────────────────
\hypersetup{
    colorlinks=true,
    linkcolor=black,
    urlcolor=black,
    citecolor=black
}

% ── Section formatting (matches PDF style) ───────────────────────────────────
\titleformat{\section}{\large\bfseries}{\thesection}{1em}{}
\titleformat{\subsection}{\normalsize\bfseries}{\thesubsection}{1em}{}
\titleformat{\subsubsection}{\normalsize\bfseries}{\thesubsubsection}{1em}{}

% ── Header / Footer ──────────────────────────────────────────────────────────
\pagestyle{fancy}
\fancyhf{}
\rhead{\small Neural Networks -- Homework 1}
\lhead{\small Cairo University}
\cfoot{\thepage}
\renewcommand{\headrulewidth}{0.4pt}

% ── Graphics path ────────────────────────────────────────────────────────────
\graphicspath{{Figures/}}

% ── Useful macros ────────────────────────────────────────────────────────────
\newcommand{\yhat}{\hat{y}}
\newcommand{\bbeta}{\boldsymbol{\beta}}

% ════════════════════════════════════════════════════════════════════════════
\begin{document}

% ── Title Page ───────────────────────────────────────────────────────────────
\begin{titlepage}
\centering
\vspace*{0.5cm}


\vspace{0.5cm}
{\Large\bfseries Cairo University\\}
\vspace{0.2cm}
{\large Faculty of Engineering\\}
\vspace{0.2cm}
{\normalsize Electronics and Electrical Communications Engineering\\}

\vspace{1.2cm}
\rule{0.8\textwidth}{1.5pt}\\[0.4cm]
{\LARGE\bfseries Homework 1}\\[0.3cm]
{\large\bfseries ELC 4028 -- Artificial Neural Networks}\\[0.2cm]
\rule{0.8\textwidth}{1.5pt}

\vspace{1.2cm}

\begin{tabular}{|c|c|l|}
\hline
\textbf{ID} & \textbf{Section} & \textbf{Name} \\
\hline
9220984 & 4 & Yousef Khaled Omar Mahmoud \\
\hline
\end{tabular}

\vspace{1.2cm}

\textit{Supervised by:}\\[0.3cm]
\textit{Dr.\ Mohsen Rashwan}\\
\textit{Dr.\ Mohamed Motaz}

\vfill
\end{titlepage}

% ── Table of Contents ────────────────────────────────────────────────────────
\setcounter{tocdepth}{3}
\tableofcontents
\newpage


% =========================
% PERCEPTRON LEARNING RULE
% =========================
\section{Perceptron Learning Rule Derivation}

In this section, we derive the perceptron learning rule step by step starting from the model definition.

% =========================
\subsection{Perceptron Model}

The perceptron computes a linear function of the input:

\[
a = \mathbf{w}^T \mathbf{x} + b
\]

The predicted label is:

\[
\hat{y} = \text{sign}(a)
\]

\textbf{Where:}
\begin{itemize}
    \item $\mathbf{x}$: input feature vector
    \item $\mathbf{w}$: weight vector
    \item $b$: bias term
    \item $a$: activation
    \item $\hat{y}$: predicted label
    \item $y \in \{+1, -1\}$: true label
\end{itemize}

% =========================
\subsection{Correct vs Incorrect Prediction}

The prediction is correct if:

\[
y(\mathbf{w}^T \mathbf{x} + b) > 0
\]

and incorrect if:

\[
y(\mathbf{w}^T \mathbf{x} + b) \leq 0
\]

\textbf{Explanation:}

\begin{itemize}
    \item If $y = +1$, we need $(\mathbf{w}^T \mathbf{x} + b) > 0$
    \item If $y = -1$, we need $(\mathbf{w}^T \mathbf{x} + b) < 0$
\end{itemize}

Multiplying by $y$ ensures the expression is positive only when the prediction is correct.

% =========================
\subsection{Objective}

If a prediction is incorrect, we aim to:
\begin{itemize}
    \item Update $\mathbf{w}$ and $b$
    \item Move the decision boundary toward correct classification
\end{itemize}

% =========================
\subsection{Proposed Update Rule}

We propose the following update:

\[
\mathbf{w}_{new} = \mathbf{w} + \eta y \mathbf{x}
\]

\[
b_{new} = b + \eta y
\]

where $\eta > 0$ is the learning rate.

\textbf{Update condition:}

\[
y(\mathbf{w}^T \mathbf{x} + b) \leq 0
\]

% =========================
\subsection{Why the Update Works}

We analyze the quantity:

\[
y(\mathbf{w}_{new}^T \mathbf{x} + b_{new})
\]

Substitute the updated parameters:

\[
= y \left( (\mathbf{w} + \eta y \mathbf{x})^T \mathbf{x} + (b + \eta y) \right)
\]

\[
= y \left( \mathbf{w}^T \mathbf{x} + \eta y \mathbf{x}^T \mathbf{x} + b + \eta y \right)
\]

\[
= y(\mathbf{w}^T \mathbf{x} + b) + \eta y^2 (\mathbf{x}^T \mathbf{x} + 1)
\]

Since:

\[
y^2 = 1
\]

we obtain:

\[
= y(\mathbf{w}^T \mathbf{x} + b) + \eta (\|\mathbf{x}\|^2 + 1)
\]

\textbf{Observation:}

\begin{itemize}
    \item $\|\mathbf{x}\|^2 + 1 > 0$
    \item Therefore, the value strictly increases after the update
\end{itemize}

\textbf{Conclusion:}

The update improves the classification score and moves the prediction toward the correct class.

% =========================
\subsection{Final Learning Rule}

\textbf{Weight update:}
\[
\mathbf{w} \leftarrow \mathbf{w} + \eta y \mathbf{x}
\]

\textbf{Bias update:}
\[
b \leftarrow b + \eta y
\]

\textbf{Apply only when:}
\[
y(\mathbf{w}^T \mathbf{x} + b) \leq 0
\]

% =========================
\subsection{Interpretation}

The perceptron learning rule:
\begin{itemize}
    \item Updates parameters only when an error occurs
    \item Moves the decision boundary in the correct direction
    \item Reinforces correct classification over time
\end{itemize}

This rule forms the foundation of linear classifiers and connects directly to margin-based methods such as Support Vector Machines.

\end{document}
